\begin{proposition}[Proved: no finite set of exterior generators in modular rank three]
\label{C:BO3_not_finitely_generated}
For every algebraically closed field $k$ of characteristic $p>2$,
the split ring $WG_{\mathrm{sp}}(BO_3)$ is not finitely generated
as a ring. More strongly, no finite subset generates it under
addition, additive inverses, multiplication and all the exterior
operations $\lambda^j$. This assertion includes
$k=\overline{\mathbb F}_p$; it does not rest on uncountability.
\end{proposition}
\begin{proof}
Use the wild Schur algebra $S_k(2,d)$ and fully faithful functor
$F$ of Proposition~\ref{C:BO3_wild_gate}. The definition of wild
type in \cite[Section 2.1, printed p.142]{C-DotyErdmannMartinNakano1999}
provides a finite-rank right-free
$S_k(2,d)$--$k\langle X,Y\rangle$ bimodule $Q$ for which
\[
 E(V)=Q\otimes_{k\langle X,Y\rangle}V
\]
preserves indecomposability and distinguishes isomorphism classes.
Choose a right basis of $Q$. The finitely many matrices describing
a basis of the finite-dimensional algebra $S_k(2,d)$ have entries
which are finite polynomials in $X,Y$ with coefficients in $k$.
Thus this family construction is defined over a finitely generated
subfield of $k$.

Suppose a finite list of classes generated the ring with its
exterior operations. Each class is a difference of finite-dimensional
symmetric representations. Their coaction matrices and forms involve
finitely many coefficients. Let $L\subset k$ be a field finitely
generated over $\mathbb F_p$ containing all these coefficients
and the coefficients of the chosen bimodule $Q$.
Every automorphism of $k$ fixing $L$ fixes each proposed generator.
It commutes with tensor products, direct sums, duals and exterior
powers, so it fixes every class generated from that finite list.

For $a\in k$ let $V_a$ be the one-dimensional
$k\langle X,Y\rangle$-module on which $X$ acts as $a$ and $Y$
as zero. These are pairwise nonisomorphic indecomposables. Hence
$M_a=F(E(V_a))$ are pairwise nonisomorphic indecomposable
$O_3$-representations. By (C.3.17), each hyperbolic class
$[H(M_a)]$ is shared by at most two parameters $a$: its only
possible ambiguity is replacing $M_a$ by its dual.
Because the construction is defined over $L$,
\[
 \sigma[H(M_a)]=[H(M_{\sigma(a)})]
 \qquad (\sigma\in\operatorname{Aut}(k/L)).             \tag{C.3.18}
\]

There is an $a\in k$ whose orbit under this automorphism group
has more than two elements. Here is an explicit field-theoretic
justification. The algebraic closure of $\mathbb F_p$ inside the
finitely generated field $L$ is a finite field $\mathbb F_{p^c}$.
Indeed a transcendence basis makes $L$ finite over a purely
transcendental field, so the degrees of its algebraic constants
are bounded; their union is therefore finite. Choose a prime
$\ell>2$ not dividing $c$ and choose
$a\in\mathbb F_{p^\ell}\setminus\mathbb F_p\subset k$.
Then $L\cap\mathbb F_{p^\ell}=\mathbb F_p$, so
$[L(a):L]=\ell$ and this extension is Galois. Its automorphisms
extend to $k$: first extend to an algebraic closure of $L$ in
$k$, then fix a transcendence basis and extend algebraically.
Thus the orbit of $a$ has size $\ell$.

At most two parameters in that orbit can have the same hyperbolic
class as $a$. Equation (C.3.18) therefore supplies an automorphism
fixing all proposed generators but moving $[H(M_a)]$, a contradiction.
This proves both non-finite-generation assertions. It is an
algebraic restriction on any future explicit description of the
modular ring, not a classification of its indecomposable generators.
\end{proof}
